% --- Archon physics formalization source begin ---
% archon:physics
% archon:covers PhyXMiniProblems/problem_phyx_mini_0875.lean
% archon:source-report reports/phyx_mini/problem_phyx_mini_0875.source.json

\chapter{Physics problem phyx\_mini\_0875}
\label{ch:PhyXMiniProblems_problem_phyx_mini_0875}

\paragraph{Problem source.}
\#\# Physical scenario

The figure shows a graph of \textbackslash{}( V \textbackslash{}) versus \textbackslash{}( x \textbackslash{}) in a region of space. The potential is independent of \textbackslash{}( y \textbackslash{}) and \textbackslash{}( z \textbackslash{}).

\#\# Question

What is \textbackslash{}( E\_\{x\} \textbackslash{}) at \textbackslash{}( x = 2\textbackslash{} \textbackslash{}mathrm\{cm\} \textbackslash{})?

\#\# Answer choices

- A: A: 12V/m

- B: B: -14V/m

- C: C: -200V/m

- D: D: -5V/m

\#\# Figure caption (auxiliary; use the image as primary evidence)

The image is a graph with two axes. The x-axis is labeled "x (m)" and the y-axis is labeled "V (V)." The graph shows a piecewise linear plot with the following characteristics:

1. The plot starts at point (-3, 0) on the x-axis and descends linearly to (-1, -10).
2. It then ascends linearly to (1, 10).
3. From (1, 10), it descends linearly to (3, 0).

The x-axis is marked at intervals of 1 unit, from -3 to 3, and the y-axis is marked from -10 to 10 in increments of 5 units. The lines connecting the points are shown in green.

\#\# Dataset metadata

Electrostatics; Physical Model Grounding Reasoning, Spatial Relation Reasoning

\paragraph{Recorded answer/context.}
D: D: -5V/m

\paragraph{Figure/image path.}
/root/proposal\_for\_physic/science-mango/phyx\_mini\_run/phyx\_data/test\_image/875.png

\paragraph{Formalization target.}
create a compiling Lean file with sorry bodies at `PhyXMiniProblems/problem\_phyx\_mini\_0875.lean`. The Lean declarations must preserve the physical quantities, dimensions or dimensional roles, figure labels, governing-law hypotheses, and final relation expressed by this problem.
Use Mathlib/Physlib names found through LeanExplore where available. If a physics API is missing, introduce faithful local abstractions rather than scalar placeholder aliases.

\begin{theorem}[Physics formalization target]
\label{thm:physics:phyx_mini_0875:target}
\lean{PhyXMiniProblems.ProblemPhyXMini0875.problem_phyx_mini_0875}
\uses{def:physics:phyx-mini-0875:phyxminiproblems-problemphyxmini0875-electrostaticpotentialgraphsetup, def:physics:phyx-mini-0875:phyxminiproblems-problemphyxmini0875-electricfieldxinvoltspermeter, def:physics:phyx-mini-0875:phyxminiproblems-problemphyxmini0875-matchessuppliedpotentialgraph, def:physics:phyx-mini-0875:phyxminiproblems-problemphyxmini0875-potentialprofilematchesfigure, def:physics:phyx-mini-0875:phyxminiproblems-problemphyxmini0875-potentialisindependentofyandz, def:physics:phyx-mini-0875:phyxminiproblems-problemphyxmini0875-calibrateselectricfieldxcomponent, def:physics:phyx-mini-0875:phyxminiproblems-problemphyxmini0875-satisfieselectrostaticpotentialfieldlaw, def:physics:phyx-mini-0875:phyxminiproblems-problemphyxmini0875-questionposition, def:physics:phyx-mini-0875:phyxminiproblems-problemphyxmini0875-answermatcheselectricfield}
The assigned autoformalize agent should translate this physics problem into Lean declarations in the covered file, with theorem and lemma proof bodies written as `by sorry`.
\end{theorem}
\begin{proof}
This is an autoformalization task, not a proof task. Produce statements that can later be proved by the physics prover without weakening the physical model.
\end{proof}

% --- Archon declaration topology begin ---
\section{Declaration topology}
\begin{definition}[electric Potential Dimension]
  \label{def:physics:phyx-mini-0875:phyxminiproblems-problemphyxmini0875-electricpotentialdimension}
  \lean{PhyXMiniProblems.ProblemPhyXMini0875.electricPotentialDimension}
  The dimension M L² T⁻² C⁻¹ of electric potential
\end{definition}
\begin{proof}
  This is the declared model definition of electric Potential Dimension.
\end{proof}
\begin{definition}[electric Field Component Dimension]
  \label{def:physics:phyx-mini-0875:phyxminiproblems-problemphyxmini0875-electricfieldcomponentdimension}
  \lean{PhyXMiniProblems.ProblemPhyXMini0875.electricFieldComponentDimension}
  The dimension M L T⁻² C⁻¹ of an electric-field component
\end{definition}
\begin{proof}
  This is the declared model definition of electric Field Component Dimension.
\end{proof}
\begin{definition}[Axial Position Quantity]
  \label{def:physics:phyx-mini-0875:phyxminiproblems-problemphyxmini0875-axialpositionquantity}
  \lean{PhyXMiniProblems.ProblemPhyXMini0875.AxialPositionQuantity}
  A signed physical position on the plotted Cartesian axis
\end{definition}
\begin{proof}
  This is the declared model definition of Axial Position Quantity.
\end{proof}
\begin{definition}[Electric Potential Quantity]
  \label{def:physics:phyx-mini-0875:phyxminiproblems-problemphyxmini0875-electricpotentialquantity}
  \lean{PhyXMiniProblems.ProblemPhyXMini0875.ElectricPotentialQuantity}
  \uses{def:physics:phyx-mini-0875:phyxminiproblems-problemphyxmini0875-electricpotentialdimension}
  A signed, unit-independent electric potential
\end{definition}
\begin{proof}
  This is the declared model definition of Electric Potential Quantity.
\end{proof}
\begin{definition}[Electric Field Component Quantity]
  \label{def:physics:phyx-mini-0875:phyxminiproblems-problemphyxmini0875-electricfieldcomponentquantity}
  \lean{PhyXMiniProblems.ProblemPhyXMini0875.ElectricFieldComponentQuantity}
  \uses{def:physics:phyx-mini-0875:phyxminiproblems-problemphyxmini0875-electricfieldcomponentdimension}
  A signed, unit-independent Cartesian electric-field component
\end{definition}
\begin{proof}
  This is the declared model definition of Electric Field Component Quantity.
\end{proof}
\begin{definition}[position Readout]
  \label{def:physics:phyx-mini-0875:phyxminiproblems-problemphyxmini0875-positionreadout}
  \lean{PhyXMiniProblems.ProblemPhyXMini0875.positionReadout}
  \uses{def:physics:phyx-mini-0875:phyxminiproblems-problemphyxmini0875-axialpositionquantity}
  Read a signed axial position in a selected physical length unit
\end{definition}
\begin{proof}
  This is the declared model definition of position Readout.
\end{proof}
\begin{definition}[position From Readout]
  \label{def:physics:phyx-mini-0875:phyxminiproblems-problemphyxmini0875-positionfromreadout}
  \lean{PhyXMiniProblems.ProblemPhyXMini0875.positionFromReadout}
  \uses{def:physics:phyx-mini-0875:phyxminiproblems-problemphyxmini0875-axialpositionquantity}
  Construct a signed physical position from a readout in a selected unit
\end{definition}
\begin{proof}
  This is the declared model definition of position From Readout.
\end{proof}
\begin{definition}[electric Potential In Volts]
  \label{def:physics:phyx-mini-0875:phyxminiproblems-problemphyxmini0875-electricpotentialinvolts}
  \lean{PhyXMiniProblems.ProblemPhyXMini0875.electricPotentialInVolts}
  \uses{def:physics:phyx-mini-0875:phyxminiproblems-problemphyxmini0875-electricpotentialquantity}
  Coherent-SI readout of electric potential, in volts
\end{definition}
\begin{proof}
  This is the declared model definition of electric Potential In Volts.
\end{proof}
\begin{definition}[electric Field Component In Volts Per Meter]
  \label{def:physics:phyx-mini-0875:phyxminiproblems-problemphyxmini0875-electricfieldcomponentinvoltspermeter}
  \lean{PhyXMiniProblems.ProblemPhyXMini0875.electricFieldComponentInVoltsPerMeter}
  \uses{def:physics:phyx-mini-0875:phyxminiproblems-problemphyxmini0875-electricfieldcomponentquantity}
  Coherent-SI readout of a signed field component, in volts per metre
\end{definition}
\begin{proof}
  This is the declared model definition of electric Field Component In Volts Per Meter.
\end{proof}
\begin{definition}[Graph Axis]
  \label{def:physics:phyx-mini-0875:phyxminiproblems-problemphyxmini0875-graphaxis}
  \lean{PhyXMiniProblems.ProblemPhyXMini0875.GraphAxis}
  The horizontal and vertical axes visible in image 875.png
\end{definition}
\begin{proof}
  This is the declared model definition of Graph Axis.
\end{proof}
\begin{definition}[Graph Axis Quantity]
  \label{def:physics:phyx-mini-0875:phyxminiproblems-problemphyxmini0875-graphaxisquantity}
  \lean{PhyXMiniProblems.ProblemPhyXMini0875.GraphAxisQuantity}
  Physical quantity named by a graph axis
\end{definition}
\begin{proof}
  This is the declared model definition of Graph Axis Quantity.
\end{proof}
\begin{definition}[Graph Axis Unit Label]
  \label{def:physics:phyx-mini-0875:phyxminiproblems-problemphyxmini0875-graphaxisunitlabel}
  \lean{PhyXMiniProblems.ProblemPhyXMini0875.GraphAxisUnitLabel}
  Literal unit text printed beside a graph axis
\end{definition}
\begin{proof}
  This is the declared model definition of Graph Axis Unit Label.
\end{proof}
\begin{definition}[Potential Graph Vertex]
  \label{def:physics:phyx-mini-0875:phyxminiproblems-problemphyxmini0875-potentialgraphvertex}
  \lean{PhyXMiniProblems.ProblemPhyXMini0875.PotentialGraphVertex}
  The four vertices of the piecewise-linear green plot, from left to right
\end{definition}
\begin{proof}
  This is the declared model definition of Potential Graph Vertex.
\end{proof}
\begin{definition}[Potential Graph Segment]
  \label{def:physics:phyx-mini-0875:phyxminiproblems-problemphyxmini0875-potentialgraphsegment}
  \lean{PhyXMiniProblems.ProblemPhyXMini0875.PotentialGraphSegment}
  The three straight green segments connecting adjacent vertices
\end{definition}
\begin{proof}
  This is the declared model definition of Potential Graph Segment.
\end{proof}
\begin{definition}[expected Segment Endpoints]
  \label{def:physics:phyx-mini-0875:phyxminiproblems-problemphyxmini0875-expectedsegmentendpoints}
  \lean{PhyXMiniProblems.ProblemPhyXMini0875.expectedSegmentEndpoints}
  \uses{def:physics:phyx-mini-0875:phyxminiproblems-problemphyxmini0875-potentialgraphvertex, def:physics:phyx-mini-0875:phyxminiproblems-problemphyxmini0875-potentialgraphsegment}
  Endpoints that the supplied image assigns to each plotted segment
\end{definition}
\begin{proof}
  This is the declared model definition of expected Segment Endpoints.
\end{proof}
\begin{definition}[expected Vertex Position In Meters]
  \label{def:physics:phyx-mini-0875:phyxminiproblems-problemphyxmini0875-expectedvertexpositioninmeters}
  \lean{PhyXMiniProblems.ProblemPhyXMini0875.expectedVertexPositionInMeters}
  \uses{def:physics:phyx-mini-0875:phyxminiproblems-problemphyxmini0875-potentialgraphvertex}
  Metre coordinate read from the image for each plotted vertex
\end{definition}
\begin{proof}
  This is the declared model definition of expected Vertex Position In Meters.
\end{proof}
\begin{definition}[expected Vertex Potential In Volts]
  \label{def:physics:phyx-mini-0875:phyxminiproblems-problemphyxmini0875-expectedvertexpotentialinvolts}
  \lean{PhyXMiniProblems.ProblemPhyXMini0875.expectedVertexPotentialInVolts}
  \uses{def:physics:phyx-mini-0875:phyxminiproblems-problemphyxmini0875-potentialgraphvertex}
  Potential, in volts, read from the image for each plotted vertex
\end{definition}
\begin{proof}
  This is the declared model definition of expected Vertex Potential In Volts.
\end{proof}
\begin{definition}[Potential Versus Position Figure]
  \label{def:physics:phyx-mini-0875:phyxminiproblems-problemphyxmini0875-potentialversuspositionfigure}
  \lean{PhyXMiniProblems.ProblemPhyXMini0875.PotentialVersusPositionFigure}
  \uses{def:physics:phyx-mini-0875:phyxminiproblems-problemphyxmini0875-axialpositionquantity, def:physics:phyx-mini-0875:phyxminiproblems-problemphyxmini0875-electricpotentialquantity, def:physics:phyx-mini-0875:phyxminiproblems-problemphyxmini0875-graphaxis, def:physics:phyx-mini-0875:phyxminiproblems-problemphyxmini0875-graphaxisquantity, def:physics:phyx-mini-0875:phyxminiproblems-problemphyxmini0875-graphaxisunitlabel, def:physics:phyx-mini-0875:phyxminiproblems-problemphyxmini0875-potentialgraphvertex, def:physics:phyx-mini-0875:phyxminiproblems-problemphyxmini0875-potentialgraphsegment}
  Typed presentation data transcribed from the supplied graph
\end{definition}
\begin{proof}
  This is the declared model definition of Potential Versus Position Figure.
\end{proof}
\begin{definition}[x Axis]
  \label{def:physics:phyx-mini-0875:phyxminiproblems-problemphyxmini0875-xaxis}
  \lean{PhyXMiniProblems.ProblemPhyXMini0875.xAxis}
  The coordinate index denoted by x in the Cartesian frame
\end{definition}
\begin{proof}
  This is the declared model definition of x Axis.
\end{proof}
\begin{definition}[point On XAxis In Meters]
  \label{def:physics:phyx-mini-0875:phyxminiproblems-problemphyxmini0875-pointonxaxisinmeters}
  \lean{PhyXMiniProblems.ProblemPhyXMini0875.pointOnXAxisInMeters}
  \uses{def:physics:phyx-mini-0875:phyxminiproblems-problemphyxmini0875-xaxis}
  A Space 3 point on the x-axis, with its coordinate read in metres
\end{definition}
\begin{proof}
  This is the declared model definition of point On XAxis In Meters.
\end{proof}
\begin{definition}[Electrostatic Potential Graph Setup]
  \label{def:physics:phyx-mini-0875:phyxminiproblems-problemphyxmini0875-electrostaticpotentialgraphsetup}
  \lean{PhyXMiniProblems.ProblemPhyXMini0875.ElectrostaticPotentialGraphSetup}
  \uses{def:physics:phyx-mini-0875:phyxminiproblems-problemphyxmini0875-axialpositionquantity, def:physics:phyx-mini-0875:phyxminiproblems-problemphyxmini0875-electricpotentialquantity, def:physics:phyx-mini-0875:phyxminiproblems-problemphyxmini0875-electricfieldcomponentquantity, def:physics:phyx-mini-0875:phyxminiproblems-problemphyxmini0875-potentialversuspositionfigure}
  Independent physical fields and observables in the electrostatic setup. Neither the field component at the question position nor any answer choice is stored here. Governing-law premises below relate these independent fields
\end{definition}
\begin{proof}
  This is the declared model definition of Electrostatic Potential Graph Setup.
\end{proof}
\begin{definition}[potential Profile In Volts]
  \label{def:physics:phyx-mini-0875:phyxminiproblems-problemphyxmini0875-potentialprofileinvolts}
  \lean{PhyXMiniProblems.ProblemPhyXMini0875.potentialProfileInVolts}
  \uses{def:physics:phyx-mini-0875:phyxminiproblems-problemphyxmini0875-positionfromreadout, def:physics:phyx-mini-0875:phyxminiproblems-problemphyxmini0875-electricpotentialinvolts, def:physics:phyx-mini-0875:phyxminiproblems-problemphyxmini0875-electrostaticpotentialgraphsetup}
  Volt readout of the physical potential profile as a function of metres
\end{definition}
\begin{proof}
  This is the declared model definition of potential Profile In Volts.
\end{proof}
\begin{definition}[electric Field XIn Volts Per Meter]
  \label{def:physics:phyx-mini-0875:phyxminiproblems-problemphyxmini0875-electricfieldxinvoltspermeter}
  \lean{PhyXMiniProblems.ProblemPhyXMini0875.electricFieldXInVoltsPerMeter}
  \uses{def:physics:phyx-mini-0875:phyxminiproblems-problemphyxmini0875-axialpositionquantity, def:physics:phyx-mini-0875:phyxminiproblems-problemphyxmini0875-electricfieldcomponentinvoltspermeter, def:physics:phyx-mini-0875:phyxminiproblems-problemphyxmini0875-electrostaticpotentialgraphsetup}
  Volts-per-metre readout of the independent x-component observable
\end{definition}
\begin{proof}
  This is the declared model definition of electric Field XIn Volts Per Meter.
\end{proof}
\begin{definition}[Matches Supplied Potential Graph]
  \label{def:physics:phyx-mini-0875:phyxminiproblems-problemphyxmini0875-matchessuppliedpotentialgraph}
  \lean{PhyXMiniProblems.ProblemPhyXMini0875.MatchesSuppliedPotentialGraph}
  \uses{def:physics:phyx-mini-0875:phyxminiproblems-problemphyxmini0875-positionreadout, def:physics:phyx-mini-0875:phyxminiproblems-problemphyxmini0875-electricpotentialinvolts, def:physics:phyx-mini-0875:phyxminiproblems-problemphyxmini0875-expectedsegmentendpoints, def:physics:phyx-mini-0875:phyxminiproblems-problemphyxmini0875-expectedvertexpositioninmeters, def:physics:phyx-mini-0875:phyxminiproblems-problemphyxmini0875-expectedvertexpotentialinvolts, def:physics:phyx-mini-0875:phyxminiproblems-problemphyxmini0875-electrostaticpotentialgraphsetup}
  The axis metadata, vertices, and straight-segment incidences visible in the primary image. This contains potential data only, not an electric-field answer
\end{definition}
\begin{proof}
  This is the declared model definition of Matches Supplied Potential Graph.
\end{proof}
\begin{definition}[Potential Profile Matches Figure]
  \label{def:physics:phyx-mini-0875:phyxminiproblems-problemphyxmini0875-potentialprofilematchesfigure}
  \lean{PhyXMiniProblems.ProblemPhyXMini0875.PotentialProfileMatchesFigure}
  \uses{def:physics:phyx-mini-0875:phyxminiproblems-problemphyxmini0875-electrostaticpotentialgraphsetup, def:physics:phyx-mini-0875:phyxminiproblems-problemphyxmini0875-potentialprofileinvolts}
  The physical potential profile is the piecewise-affine green graph. The three formulas are the straight-line interpolants through the four supplied vertices and contain no electric-field conclusion
\end{definition}
\begin{proof}
  This is the declared model definition of Potential Profile Matches Figure.
\end{proof}
\begin{definition}[Potential Is Independent Of YAnd Z]
  \label{def:physics:phyx-mini-0875:phyxminiproblems-problemphyxmini0875-potentialisindependentofyandz}
  \lean{PhyXMiniProblems.ProblemPhyXMini0875.PotentialIsIndependentOfYAndZ}
  \uses{def:physics:phyx-mini-0875:phyxminiproblems-problemphyxmini0875-positionfromreadout, def:physics:phyx-mini-0875:phyxminiproblems-problemphyxmini0875-xaxis, def:physics:phyx-mini-0875:phyxminiproblems-problemphyxmini0875-electrostaticpotentialgraphsetup}
  The problem statement's assertion that potential is independent of y,z
\end{definition}
\begin{proof}
  This is the declared model definition of Potential Is Independent Of YAnd Z.
\end{proof}
\begin{definition}[Calibrates Electric Field XComponent]
  \label{def:physics:phyx-mini-0875:phyxminiproblems-problemphyxmini0875-calibrateselectricfieldxcomponent}
  \lean{PhyXMiniProblems.ProblemPhyXMini0875.CalibratesElectricFieldXComponent}
  \uses{def:physics:phyx-mini-0875:phyxminiproblems-problemphyxmini0875-positionreadout, def:physics:phyx-mini-0875:phyxminiproblems-problemphyxmini0875-xaxis, def:physics:phyx-mini-0875:phyxminiproblems-problemphyxmini0875-pointonxaxisinmeters, def:physics:phyx-mini-0875:phyxminiproblems-problemphyxmini0875-electrostaticpotentialgraphsetup, def:physics:phyx-mini-0875:phyxminiproblems-problemphyxmini0875-electricfieldxinvoltspermeter}
  The dimensionful component observable calibrates the x component of Physlib's full vector electric field in coherent SI units
\end{definition}
\begin{proof}
  This is the declared model definition of Calibrates Electric Field XComponent.
\end{proof}
\begin{definition}[Satisfies Electrostatic Potential Field Law]
  \label{def:physics:phyx-mini-0875:phyxminiproblems-problemphyxmini0875-satisfieselectrostaticpotentialfieldlaw}
  \lean{PhyXMiniProblems.ProblemPhyXMini0875.SatisfiesElectrostaticPotentialFieldLaw}
  \uses{def:physics:phyx-mini-0875:phyxminiproblems-problemphyxmini0875-positionreadout, def:physics:phyx-mini-0875:phyxminiproblems-problemphyxmini0875-electrostaticpotentialgraphsetup, def:physics:phyx-mini-0875:phyxminiproblems-problemphyxmini0875-potentialprofileinvolts, def:physics:phyx-mini-0875:phyxminiproblems-problemphyxmini0875-electricfieldxinvoltspermeter}
  Electrostatic governing law: the field is static and, wherever the plotted potential has derivative slope in volts per metre, Eₓ = -slope. This is a general physical law rather than the requested numerical answer
\end{definition}
\begin{proof}
  This is the declared model definition of Satisfies Electrostatic Potential Field Law.
\end{proof}
\begin{definition}[question Position]
  \label{def:physics:phyx-mini-0875:phyxminiproblems-problemphyxmini0875-questionposition}
  \lean{PhyXMiniProblems.ProblemPhyXMini0875.questionPosition}
  \uses{def:physics:phyx-mini-0875:phyxminiproblems-problemphyxmini0875-axialpositionquantity, def:physics:phyx-mini-0875:phyxminiproblems-problemphyxmini0875-positionfromreadout}
  The position written in the question: x = 2 cm
\end{definition}
\begin{proof}
  This is the declared model definition of question Position.
\end{proof}
\begin{definition}[Answer Choice]
  \label{def:physics:phyx-mini-0875:phyxminiproblems-problemphyxmini0875-answerchoice}
  \lean{PhyXMiniProblems.ProblemPhyXMini0875.AnswerChoice}
  Labels of the four displayed answer choices
\end{definition}
\begin{proof}
  This is the declared model definition of Answer Choice.
\end{proof}
\begin{definition}[field Component In Volts Per Meter]
  \label{def:physics:phyx-mini-0875:phyxminiproblems-problemphyxmini0875-answerchoice-fieldcomponentinvoltspermeter}
  \lean{PhyXMiniProblems.ProblemPhyXMini0875.AnswerChoice.fieldComponentInVoltsPerMeter}
  \uses{def:physics:phyx-mini-0875:phyxminiproblems-problemphyxmini0875-answerchoice}
  Signed field component, in volts per metre, printed beside each choice
\end{definition}
\begin{proof}
  This is the declared model definition of field Component In Volts Per Meter.
\end{proof}
\begin{definition}[recorded Dataset Answer]
  \label{def:physics:phyx-mini-0875:phyxminiproblems-problemphyxmini0875-recordeddatasetanswer}
  \lean{PhyXMiniProblems.ProblemPhyXMini0875.recordedDatasetAnswer}
  \uses{def:physics:phyx-mini-0875:phyxminiproblems-problemphyxmini0875-answerchoice}
  Answer label recorded by the supplied dataset metadata
\end{definition}
\begin{proof}
  This is the declared model definition of recorded Dataset Answer.
\end{proof}
\begin{definition}[Answer Matches Electric Field]
  \label{def:physics:phyx-mini-0875:phyxminiproblems-problemphyxmini0875-answermatcheselectricfield}
  \lean{PhyXMiniProblems.ProblemPhyXMini0875.AnswerMatchesElectricField}
  \uses{def:physics:phyx-mini-0875:phyxminiproblems-problemphyxmini0875-electrostaticpotentialgraphsetup, def:physics:phyx-mini-0875:phyxminiproblems-problemphyxmini0875-electricfieldxinvoltspermeter, def:physics:phyx-mini-0875:phyxminiproblems-problemphyxmini0875-questionposition, def:physics:phyx-mini-0875:phyxminiproblems-problemphyxmini0875-answerchoice}
  A choice matches the independently modeled field-component observable
\end{definition}
\begin{proof}
  This is the declared model definition of Answer Matches Electric Field.
\end{proof}
\begin{lemma}[question Position readout centimeters]
  \label{lem:physics:phyx-mini-0875:phyxminiproblems-problemphyxmini0875-questionposition-readout-centimeters}
  \lean{PhyXMiniProblems.ProblemPhyXMini0875.questionPosition_readout_centimeters}
  \uses{def:physics:phyx-mini-0875:phyxminiproblems-problemphyxmini0875-positionreadout, def:physics:phyx-mini-0875:phyxminiproblems-problemphyxmini0875-questionposition}
  The literal question coordinate has the readout 2 cm
\end{lemma}
\begin{proof}
  The conclusion follows from the governing relations and figure readouts named in the statement of question Position readout centimeters.
\end{proof}
\begin{lemma}[question Position readout meters]
  \label{lem:physics:phyx-mini-0875:phyxminiproblems-problemphyxmini0875-questionposition-readout-meters}
  \lean{PhyXMiniProblems.ProblemPhyXMini0875.questionPosition_readout_meters}
  \uses{def:physics:phyx-mini-0875:phyxminiproblems-problemphyxmini0875-positionreadout, def:physics:phyx-mini-0875:phyxminiproblems-problemphyxmini0875-questionposition}
  Equivalently, the literal question coordinate is 0.02 m
\end{lemma}
\begin{proof}
  The conclusion follows from the governing relations and figure readouts named in the statement of question Position readout meters.
\end{proof}
% --- Archon declaration topology end ---
% --- Archon physics formalization source end ---
